
\documentclass[11pt,a4paper]{article}

\usepackage[utf8]{inputenc}
\usepackage[T1]{fontenc}
\usepackage[english]{babel}
\usepackage{lmodern}

\usepackage[margin=2.5cm]{geometry}
\usepackage{booktabs}         
\midrule, \bottomrule)
\usepackage{longtable}         
\usepackage{array}
\usepackage{tabularx}
\usepackage{xcolor}
\usepackage{graphicx}
\usepackage{amsmath}
\usepackage{enumitem}
\usepackage{fancyhdr}
\usepackage[colorlinks=true,linkcolor=black,urlcolor=blue,citecolor=blue]{hyperref}

% --- Colours -----------------------------------------------------------------
\definecolor{accent}{HTML}{1E3A5F}
\definecolor{softgrey}{HTML}{F5F5F5}
\definecolor{warnred}{HTML}{C0392B}
\definecolor{okgreen}{HTML}{1E8449}

% --- Header / footer ---------------------------------------------------------
\pagestyle{fancy}
\fancyhf{}
\lhead{\small Bias and Fairness Analysis}
\rhead{\small RSF --- AML Overall Survival}
\cfoot{\thepage}
\renewcommand{\headrulewidth}{0.4pt}

% --- Handy macros ------------------------------------------------------------
\newcommand{\var}[1]{\texttt{#1}}                      % variable names
\newcommand{\bad}[1]{\textcolor{warnred}{\textbf{#1}}} % highlight a problem
\newcommand{\good}[1]{\textcolor{okgreen}{\textbf{#1}}}% highlight a good result

% Boxed key-finding paragraph
\usepackage{mdframed}
\newmdenv[
  backgroundcolor=softgrey,
  linecolor=accent,
  linewidth=2pt,
  topline=false, bottomline=false, rightline=false,
  innerleftmargin=10pt, innertopmargin=8pt, innerbottommargin=8pt,
  skipabove=10pt, skipbelow=10pt
]{keybox}

% =============================================================================
\title{\vspace{-1.5cm}\textbf{Bias and Fairness Analysis}\\[0.3cm]
       \large Overall Survival in Acute Myeloid Leukemia}

\date{\today}

\begin{document}
\maketitle
\thispagestyle{fancy}

\begin{abstract}
\noindent
A global C-index of 0.752 describes a model's average behaviour but says nothing about
\emph{who} it serves well. This report audits a Random Survival Forest trained on 3{,}323 AML (Acute Myeloid Leukemia)
patients across 10 European centres, examining whether predictive quality is distributed
equitably across patient subgroups. Two independent methods establish that the model does not
rely on the treating hospital. However, calibration analysis reveals a systematic and
reproducible optimism bias for adverse-risk cytogenetic groups, invisible to discrimination
metrics.
\end{abstract}
\newpage
\tableofcontents
\newpage


\section{Introduction}


A global performance figure describes a model's \emph{average} behaviour. It says nothing about
\textbf{who} the model serves well and who it serves poorly. A model reaching a C-index of
0.752 overall may be excellent for the majority of patients and unreliable for a specific 
and potentially vulnerable subgroup.

Bias analysis addresses a question that aggregate metrics cannot: \textbf{is predictive quality
distributed equitably across patient subgroups, or does the model systematically disadvantage
an identifiable population?}

In a clinical context, this question is not merely statistical. A prognostic tool that
overestimates survival for high-risk patients could contribute to under-treatment, delayed
referral to intensive care pathways, or inappropriate reassurance. \emph{The direction of an
error matters as much as its magnitude.}

% =============================================================================
\section{A taxonomy of bias}
% =============================================================================

Bias in a machine learning pipeline can originate at four distinct stages. Distinguishing them
matters because each requires a different remedy.

\subsection{Data-level bias}

Introduced before any model is trained.

\begin{longtable}{@{}p{3.6cm} p{5.5cm} p{5.5cm}@{}}
\toprule
\textbf{Type} & \textbf{Definition} & \textbf{Typical consequence} \\
\midrule
\endhead
Sampling / representation & Some subgroups are under-represented relative to the target
population & The model learns those groups poorly \\
\addlinespace
Selection & Inclusion criteria systematically exclude a category of patients & Predictions do
not generalise beyond the recruited population \\
\addlinespace
Survivorship & Individuals who experienced the outcome early are absent by construction &
Survival is structurally overestimated \\
\addlinespace
Measurement & A variable is recorded differently across sources or periods & The model learns
measurement artefacts rather than biology \\
\addlinespace
Annotation / documentation & Missing data is concentrated in specific groups rather than random
& Missingness itself becomes an informative and unfair signal \\
\bottomrule
\end{longtable}

\subsection{Model-level bias}

Arises from the learning algorithm itself, even on well-balanced data.

\begin{longtable}{@{}p{3.6cm} p{5.5cm} p{5.5cm}@{}}
\toprule
\textbf{Type} & \textbf{Definition} & \textbf{Typical consequence} \\
\midrule
\endhead
Calibration & Predicted probabilities deviate systematically from observed frequencies &
Correct ranking, but wrong absolute risk \\
\addlinespace
Regression to the mean & Ensemble averaging pulls extreme predictions toward the distribution
centre & Severe and mild cases are both under-differentiated \\
\addlinespace
Discrimination & Ranking quality differs across subgroups & The model is less useful for some
patients \\
\addlinespace
Aggregation & A single model is fitted where subgroups follow different underlying
relationships & Neither subgroup is optimally served \\
\bottomrule
\end{longtable}

\subsection{Leakage and confounding}

\begin{longtable}{@{}p{3.6cm} p{5.5cm} p{5.5cm}@{}}
\toprule
\textbf{Type} & \textbf{Definition} & \textbf{Typical consequence} \\
\midrule
\endhead
Contextual leakage & A technical or administrative variable carries outcome signal &
Performance does not transfer to a new setting \\
\addlinespace
Institutional & The model relies on the treating centre rather than patient characteristics &
Prognosis depends on place of care, not biology \\
\addlinespace
Proxy & A neutral variable encodes a protected or structural attribute & Indirect
discrimination \\
\bottomrule
\end{longtable}

\subsection{Deployment bias}

\begin{longtable}{@{}p{3.6cm} p{5.5cm} p{5.5cm}@{}}
\toprule
\textbf{Type} & \textbf{Definition} & \textbf{Typical consequence} \\
\midrule
\endhead
Distribution shift & The deployment population differs from the training population & Silent
degradation of performance \\
\addlinespace
Out-of-distribution input & The interface allows inputs never observed in training &
Predictions produced with no empirical basis \\
\addlinespace
Automation & Users over-trust model output & Amplification of any residual bias \\
\bottomrule
\end{longtable}

% =============================================================================
\section{Why survival models require a specific fairness framework}
% =============================================================================

Standard fairness metrics demographic parity, equalised odds, equal opportunity  assume
a \textbf{classification} setting with fully observed binary outcomes. Neither assumption holds
here.

\paragraph{The outcome is a duration, not a label.}
There is no confusion matrix, therefore no false positive or false negative rate to equalise
across groups.

\paragraph{The outcome is right-censored for approximately 48\% of patients.}
For nearly half the cohort we know only that the patient was alive at last follow-up. This has
two methodological consequences:

\begin{itemize}[leftmargin=*]
  \item \textbf{The C-index depends on the censoring distribution.} A subgroup with short
        follow-up produces fewer comparable pairs, and therefore an unstable estimate. Raw
        differences between groups cannot be interpreted as bias without controlling for this.
  \item \textbf{Kaplan--Meier is the only unbiased reference} for observed survival, and its
        reliability degrades as censoring increases.
\end{itemize}

\paragraph{Discrimination and calibration are independent failure modes.}
A model may rank patients correctly within a subgroup while being systematically wrong about
their absolute survival. Clinically, the second error is the more dangerous, and it is entirely
invisible to the C-index. Any credible fairness audit of a survival model must therefore
evaluate \textbf{both}.

% =============================================================================
\section{Application to our study}
% =============================================================================

\subsection{Available axes}

The dataset contains: \var{ID}, \var{CENTER}, \var{BM\_BLAST}, \var{WBC}, \var{ANC},
\var{MONOCYTES}, \var{HB}, \var{PLT}, \var{CYTOGENETICS}, \var{CYTO\_CLASS}, \var{OS\_YEARS},
\var{OS\_STATUS}, plus patient-level molecular aggregates.

\begin{keybox}
\textbf{No age and no sex variable is available.} Demographic fairness --- the axis regulators
and reviewers examine first --- therefore \textbf{cannot be assessed}. This is stated as a
limitation rather than worked around. It is also a modelling weakness in its own right: age is
the second strongest prognostic factor in AML after cytogenetics.
\end{keybox}

Three axes remain analysable:

\begin{table}[h]
\centering
\begin{tabular}{@{}p{2.8cm} p{3.5cm} p{7.5cm}@{}}
\toprule
\textbf{Axis} & \textbf{Levels} & \textbf{Fairness question} \\
\midrule
\var{CENTER} & 23 hospital centres & Does prognosis depend on where the patient was treated? \\
\var{CYTO\_CLASS} & 8 ELN 2022 risk groups & Is the model as reliable for adverse-risk as for
favourable-risk patients? \\
\var{HAS\_MOL\_DATA} & tested / not tested & Does the model underserve patients without
molecular profiling? \\
\bottomrule
\end{tabular}
\end{table}

\var{CENTER} is the most sensitive axis, because it is \textbf{a feature of the model itself}.
The RSF could have learned institutional differences in care, recruitment or documentation
rather than patient biology, a textbook case of institutional bias.

\subsection{Confounders identified before analysis}

\paragraph{Censoring is highly heterogeneous across centres}, ranging from 8.4\% to 90.2\%, and
correlates with apparent survival:

\begin{table}[h]
\centering
\begin{tabular}{@{}lrrrr@{}}
\toprule
\textbf{Centre} & \textbf{n} & \textbf{Censored} & \textbf{KM $S(2y)$} & \textbf{Max follow-up} \\
\midrule
MUV  &  83 &  8.4\% & 0.446 & 12.79 y \\
UOB  &  86 & 18.6\% & 0.581 & 14.70 y \\
CCH  & 140 & 76.4\% & 0.834 &  9.32 y \\
ROM  & 102 & 90.2\% & 0.809 &  2.84 y \\
\bottomrule
\end{tabular}
\end{table}

Centres with short follow-up appear to have better outcomes simply because fewer deaths were
observed. This is a censoring artefact, not a clinical difference.

\paragraph{Subgroup sizes are severely imbalanced}, from 786 patients (KI) to 6 (REL), and from
1{,}780 (\var{Normal}) to 79 (\var{Complex}). Estimates below roughly 50 patients are
statistically uninterpretable.

% =============================================================================
\section{Methodology}
% =============================================================================

\subsection{Discrimination per subgroup}

C-index computed \textbf{within} each subgroup, with \textbf{bootstrap 95\% confidence
intervals} (200 resamples). Subgroups below 50 patients or 15 events are reported but not
interpreted.

\begin{keybox}
\textbf{Decision rule:} bias is concluded only where confidence intervals are
\textbf{disjoint} never from raw gaps. With subgroup sizes ranging over two orders of
magnitude, overlapping intervals are the expected outcome even under perfect fairness.
\end{keybox}

\subsection{Calibration per subgroup}

Mean predicted survival at 2 years compared against the \textbf{Kaplan--Meier} estimate
observed in the same subgroup:

\begin{equation}
\text{gap} \;=\; \overline{\hat{S}(2\text{y})} \;-\; S_{\text{KM}}(2\text{y})
\end{equation}

A \textbf{positive gap means the model is too optimistic} for that group. Censoring rate is
reported alongside each gap, since a heavily censored subgroup yields an unreliable KM
reference.

\subsection{Dependence on institutional variables}

\textbf{Permutation test:} the \var{CENTER} column is randomly shuffled and the loss in C-index
measured (5 repetitions). If the model genuinely relies on the treating hospital, destroying
that information should degrade performance. A drop below $\sim$0.005 indicates negligible
dependence. This is cross-checked against the SHAP analysis conducted independently at the
individual patient level.

\subsection{Overfitting control}

Because subgroup metrics are computed on training data, the calibration analysis is repeated on
a held-out split (25\%, stratified on event status) to verify that any observed miscalibration
is not an artefact of memorisation.

% =============================================================================
\section{Results}
% =============================================================================

\subsection{Discrimination --- no detectable bias}

\begin{table}[h]
\centering
\begin{tabular}{@{}lccl@{}}
\toprule
\textbf{Axis} & \textbf{Observed range} & \textbf{Evaluable} & \textbf{Conclusion} \\
\midrule
\var{CENTER}        & 0.752 -- 0.896 & 11 / 23 & all CIs overlap $\rightarrow$ no bias \\
\var{CYTO\_CLASS}   & 0.743 -- 0.817 &  7 / 8  & weak, non-significant trend \\
\var{HAS\_MOL\_DATA}& 0.798 vs 0.810 &  2 / 2  & no access-related inequity \\
\bottomrule
\end{tabular}
\end{table}

Selected centre results:

\begin{table}[h]
\centering
\begin{tabular}{@{}lrrcc@{}}
\toprule
\textbf{Centre} & \textbf{n} & \textbf{events} & \textbf{C-index} & \textbf{95\% CI} \\
\midrule
CGM  & 107 &  37 & 0.896 & [0.848 -- 0.932] \\
PV   & 316 &  81 & 0.861 & [0.828 -- 0.898] \\
DUS  & 445 & 288 & 0.807 & [0.784 -- 0.832] \\
KI   & 786 & 483 & 0.800 & [0.781 -- 0.820] \\
RMCN & 182 & 105 & 0.752 & [0.707 -- 0.802] \\
\bottomrule
\end{tabular}
\end{table}

Even the extremes (CGM at 0.896, RMCN at 0.752) are not cleanly separated once uncertainty is
accounted for. \textbf{The model ranks patients with comparable effectiveness across all
evaluable subgroups.}

By cytogenetic class, a mild gradient appears, \var{Normal} / \var{5q deletion} around
0.80--0.82 against \var{Monosomy 7} (0.743) and \var{Complex} (0.747) but intervals
overlap. The most plausible interpretation is structural rather than a model defect: in
very-poor-prognosis groups most patients die quickly, leaving less to discriminate.

\subsection{Institutional dependence --- ruled out}

\begin{table}[h]
\centering
\begin{tabular}{@{}lr@{}}
\toprule
\textbf{Test} & \textbf{Result} \\
\midrule
C-index (reference)                & 0.8105 \\
Drop when \var{CENTER} is shuffled & \good{0.0051 $\pm$ 0.0002} \\
Relative impact                    & 0.6\% of performance \\
\bottomrule
\end{tabular}
\end{table}

Destroying all information about the treating hospital costs the model 0.5\% of a C-index
point. This is corroborated independently by the SHAP analysis, where \emph{Treatment center}
contributes $\approx 0$ to individual predictions.

\begin{keybox}
\good{Prognosis does not depend on where the patient was treated.} This is a negative result,
and a strong one: the variable is present in the model but effectively inert.
\end{keybox}

\subsection{Calibration --- systematic bias on adverse-risk groups}

This is the principal finding.

\begin{table}[h]
\centering
\small
\begin{tabular}{@{}lrrrrrr@{}}
\toprule
\textbf{Group} & \textbf{n} & \textbf{Predicted} & \textbf{Observed} & \textbf{Gap} &
\textbf{Held-out} & \textbf{Censored} \\
\midrule
\bad{Complex}    &   79 & 0.425 & 0.269 & \bad{+0.156} & \bad{+0.148} & 27.8\% \\
\bad{Monosomy 7} &  164 & 0.345 & 0.246 & \bad{+0.099} & \bad{+0.102} & 25.0\% \\
Trisomy 8        &  181 & 0.479 & 0.428 & +0.052 & +0.021 & 35.9\% \\
Missing          &  285 & 0.600 & 0.591 & +0.009 & +0.004 & 40.4\% \\
5q deletion      &  296 & 0.567 & 0.560 & +0.007 & $-$0.051 & 41.2\% \\
Other            &  304 & 0.599 & 0.608 & $-$0.010 & +0.016 & 48.7\% \\
Normal           & 1780 & 0.679 & 0.699 & $-$0.020 & $-$0.024 & 57.2\% \\
\bottomrule
\end{tabular}
\end{table}

\textbf{The model is systematically too optimistic for patients with adverse cytogenetics.} For
complex karyotypes it predicts 42.5\% survival at two years where 26.9\% is observed a
15.6-point overestimation.

Three elements establish this as a genuine model bias rather than noise:

\begin{enumerate}[leftmargin=*]
  \item \textbf{The gradient is monotonic.} The worse the prognosis, the larger the
        overestimation. Favourable and majority groups are well calibrated
        ($|\text{gap}| < 0.02$).
  \item \textbf{It reproduces out of sample.} Held-out gaps ($+0.148$, $+0.102$) match training
        gaps ($+0.156$, $+0.099$) almost exactly, excluding overfitting.
  \item \textbf{It is not a censoring artefact.} \var{Monosomy 7} has only 14.7\% censoring in
        the held-out split, so its Kaplan--Meier reference is reliable.
\end{enumerate}

\paragraph{Mechanism.}
This is \textbf{regression to the mean}, an expected property of random forests: each
prediction aggregates leaf averages, pulling extreme cases toward the centre of the outcome
distribution. The affected groups are simultaneously the most severe and the least represented
(79 and 164 patients), which amplifies the effect. The bias therefore \emph{originates in the
data} (under-representation of severe cases) and \emph{manifests in the model} (shrinkage of
their predicted risk).

\begin{keybox}
\textbf{Clinical implication.} The error runs in the most harmful direction: announcing a more
favourable prognosis than reality to the highest-risk patients. It is entirely invisible to the
C-index. \var{Complex} retains a respectable discrimination of 0.747 and surfaces only
through calibration analysis. This validates the two-pronged methodology.
\end{keybox}

% =============================================================================
\section{Bias inventory}
% =============================================================================

\subsection{Confirmed}

\begin{longtable}{@{}p{4.2cm} p{6.5cm} p{3.8cm}@{}}
\toprule
\textbf{Bias} & \textbf{Evidence} & \textbf{Origin} \\
\midrule
\endhead
Calibration / regression to the mean & $+0.156$ (Complex), $+0.099$ (Monosomy 7), reproduced
out-of-sample & Model, driven by data imbalance \\
\addlinespace
Representation & 786 $\rightarrow$ 6 patients per centre; 1{,}780 $\rightarrow$ 79 per
cytogenetic class & Data \\
\addlinespace
Censoring & 8.4\% $\rightarrow$ 90.2\% across centres, correlated with apparent survival &
Data / study design \\
\addlinespace
Documentation & 285 patients with \var{CYTO\_CLASS = Missing}, 256 without molecular testing &
Data collection \\
\addlinespace
Contextual leakage (suspected) & \var{DEPTH\_MEAN} / \var{DEPTH\_MAX} carry prognostic signal
despite being purely technical & Feature engineering \\
\bottomrule
\end{longtable}

\subsection{Investigated and ruled out}

\begin{table}[h]
\centering
\begin{tabular}{@{}p{4.5cm} p{9cm}@{}}
\toprule
\textbf{Bias} & \textbf{Evidence} \\
\midrule
Institutional & \var{CENTER} permutation costs 0.0051 C-index; SHAP contribution $\approx 0$ \\
Discrimination & All subgroup confidence intervals overlap \\
Access-to-care inequity & \var{HAS\_MOL\_DATA}: 0.798 vs 0.810, overlapping intervals \\
\bottomrule
\end{tabular}
\end{table}

\subsection{Not assessable}

\begin{table}[h]
\centering
\begin{tabular}{@{}p{4.5cm} p{9cm}@{}}
\toprule
\textbf{Bias} & \textbf{Reason} \\
\midrule
Demographic & Neither age nor sex present in the dataset \\
Selection & No information on excluded patients or referral patterns \\
Survivorship & Patients dying before karyotyping are absent by construction \\
\bottomrule
\end{tabular}
\end{table}

% =============================================================================
\section{Discussion}
% =============================================================================

The analysis produces one strong negative result and one strong positive result, and the
contrast between them is instructive.

\paragraph{The negative result is reassuring.}
Despite \var{CENTER} being an explicit model feature (a configuration that would normally
raise concern) the model does not exploit it. Two independent methods (permutation testing
at the population level, SHAP at the individual level) converge on the same conclusion. The
model bases its prognosis on biology, not on institutional context.

\paragraph{The positive result is actionable.}
The calibration bias is systematic, reproducible, and runs in the clinically harmful direction.
Critically, it would have gone entirely unnoticed had the audit been limited to discrimination
metrics: \var{Complex} shows a C-index of 0.747, which any standard evaluation would deem
acceptable.

This illustrates the central methodological point of the report: \textbf{for survival models,
ranking quality and absolute risk accuracy are separate properties, and fairness must be
assessed on both.} A model can rank a subgroup perfectly while being systematically wrong about
how long its patients will actually live.

The mechanism also clarifies where remediation should occur. Since the bias stems from
under-representation of severe cases combined with ensemble averaging, changing algorithm is
unlikely to resolve it. Post-hoc recalibration or reweighting of severe cases are the
appropriate directions.

% =============================================================================
\section{Limitations}
% =============================================================================

\begin{itemize}[leftmargin=*]
  \item \textbf{Metrics are computed on training data.} Subgroup C-indices around 0.80
        correspond to the training C-index (0.8105), not to the held-out platform score
        (0.752). Absolute values are optimistic; only \emph{relative comparisons between
        subgroups} are interpretable.
  \item \textbf{The held-out split is imperfect.} It is drawn from patients the model saw
        during training. A cross-validated calibration, measured on each validation fold, would
        be required for publication-grade evidence.
  \item \textbf{Small subgroups.} 12 of 23 centres fall below the interpretability threshold.
        Held-out \var{Complex} and \var{Monosomy 7} rest on 22 and 34 patients respectively 
        (indicative results requiring confirmation on an external cohort).
  \item \textbf{No demographic variables}, so the most commonly examined fairness axis is
        inaccessible.
  \item \textbf{A single horizon} (2 years) was used for calibration. Bias may behave
        differently at 1 or 5 years.
  \item \textbf{Out-of-distribution inputs in the application.} The Streamlit interface offers
        cytogenetic classes absent from training (\var{inv(16)}, \var{t(8;21)}) or represented
        by a single patient (\var{APL t(15;17)}), which would produce predictions with no
        empirical basis. This is a deployment bias identified during the audit and requiring
        correction before any use.
\end{itemize}

\end{document}